% !TITLE 明清
% !SUMMARY 复古与革新的张力
% !ORDER 0

\begin{intro}
明清两代近六百年，是中国古典诗歌发展的最后阶段。它不再有唐诗的气象、宋词的辉煌，却在复古与革新、模仿与突破的张力中，留下了许多值得品味的篇章。

明代的诗坛，从明初的高启、刘基，到台阁体的雍容，再到前后七子的复古运动，始终在"学唐"还是"宗宋"的争论中摇摆。前后七子提出"文必秦汉，诗必盛唐"，力图恢复唐诗的雄浑气象，但过度拟古往往导致缺乏真情实感。倒是明中叶的唐寅，不依附任何流派，以"桃花坞里桃花庵，桃花庵下桃花仙"的洒脱，写出了一种晚明文人特有的风流自赏。

清代是古典诗歌的总结期。钱谦益、吴伟业等遗民诗人以诗存史，记录了朝代更迭的剧痛。到了乾嘉年间，袁枚性灵派主张"诗者，人之性情也"，反对格调、神韵诸说的束缚，为诗坛注入了一股清新的空气。但清代诗歌最动人的，或许还是纳兰性德。他的词不作高论，只是写寻常的离别相思，"人生若只如初见，何事秋风悲画扇"——这种直逼人心的情感力量，在清代是少见的。

明清两代更大的贡献在于小说和戏曲。《牡丹亭》"原来姹紫嫣红开遍，似这般都付与断井颓垣"，《红楼梦》里借林黛玉、薛宝钗之手写下的诗词——这些作品中的诗词，已经与人物命运、情节发展融为一体，获得了前所未有的叙事深度。从这个意义上说，明清诗歌虽然不再是文学的主角，却以一种新的方式延续着它的生命。
\end{intro}
